% Options for packages loaded elsewhere
% Options for packages loaded elsewhere
\PassOptionsToPackage{unicode}{hyperref}
\PassOptionsToPackage{hyphens}{url}
\PassOptionsToPackage{dvipsnames,svgnames,x11names}{xcolor}
%
\documentclass[
  11pt,
  a4paper,
  twoside,
  openright]{scrbook}
\usepackage{xcolor}
\usepackage[inner=25mm,top=25mm,bottom=25mm,textwidth=115mm,marginparsep=7mm,marginparwidth=48mm]{geometry}
\usepackage{amsmath,amssymb}
\setcounter{secnumdepth}{5}
\usepackage{iftex}
\ifPDFTeX
  \usepackage[T1]{fontenc}
  \usepackage[utf8]{inputenc}
  \usepackage{textcomp} % provide euro and other symbols
\else % if luatex or xetex
  \usepackage{unicode-math} % this also loads fontspec
  \defaultfontfeatures{Scale=MatchLowercase}
  \defaultfontfeatures[\rmfamily]{Ligatures=TeX,Scale=1}
\fi
\usepackage{lmodern}
\ifPDFTeX\else
  % xetex/luatex font selection
  \setmainfont[Path=./fonts/,Extension=.ttf,UprightFont=*-Regular,BoldFont=*-Bold,ItalicFont=*-Italic,BoldItalicFont=*-BoldItalic]{Sarabun}
  \setsansfont[Path=./fonts/,Extension=.ttf,UprightFont=*-Regular,BoldFont=*-Bold,ItalicFont=*-Italic,BoldItalicFont=*-BoldItalic]{Sarabun}
  \setmonofont[Extension=.ttf,BoldFont=*-Bold,ItalicFont=*-Oblique,BoldItalicFont=*-BoldOblique,Scale=MatchLowercase]{DejaVuSansMono}
  \setmathfont[]{latinmodern-math.otf}
\fi
% Use upquote if available, for straight quotes in verbatim environments
\IfFileExists{upquote.sty}{\usepackage{upquote}}{}
\IfFileExists{microtype.sty}{% use microtype if available
  \usepackage[]{microtype}
  \UseMicrotypeSet[protrusion]{basicmath} % disable protrusion for tt fonts
}{}
\usepackage{setspace}
\makeatletter
\@ifundefined{KOMAClassName}{% if non-KOMA class
  \IfFileExists{parskip.sty}{%
    \usepackage{parskip}
  }{% else
    \setlength{\parindent}{0pt}
    \setlength{\parskip}{6pt plus 2pt minus 1pt}}
}{% if KOMA class
  \KOMAoptions{parskip=half}}
\makeatother
% Make \paragraph and \subparagraph free-standing
\makeatletter
\ifx\paragraph\undefined\else
  \let\oldparagraph\paragraph
  \renewcommand{\paragraph}{
    \@ifstar
      \xxxParagraphStar
      \xxxParagraphNoStar
  }
  \newcommand{\xxxParagraphStar}[1]{\oldparagraph*{#1}\mbox{}}
  \newcommand{\xxxParagraphNoStar}[1]{\oldparagraph{#1}\mbox{}}
\fi
\ifx\subparagraph\undefined\else
  \let\oldsubparagraph\subparagraph
  \renewcommand{\subparagraph}{
    \@ifstar
      \xxxSubParagraphStar
      \xxxSubParagraphNoStar
  }
  \newcommand{\xxxSubParagraphStar}[1]{\oldsubparagraph*{#1}\mbox{}}
  \newcommand{\xxxSubParagraphNoStar}[1]{\oldsubparagraph{#1}\mbox{}}
\fi
\makeatother

\usepackage{color}
\usepackage{fancyvrb}
\newcommand{\VerbBar}{|}
\newcommand{\VERB}{\Verb[commandchars=\\\{\}]}
\DefineVerbatimEnvironment{Highlighting}{Verbatim}{commandchars=\\\{\}}
% Add ',fontsize=\small' for more characters per line
\usepackage{framed}
\definecolor{shadecolor}{RGB}{241,243,245}
\newenvironment{Shaded}{\begin{snugshade}}{\end{snugshade}}
\newcommand{\AlertTok}[1]{\textcolor[rgb]{0.68,0.00,0.00}{#1}}
\newcommand{\AnnotationTok}[1]{\textcolor[rgb]{0.37,0.37,0.37}{#1}}
\newcommand{\AttributeTok}[1]{\textcolor[rgb]{0.40,0.45,0.13}{#1}}
\newcommand{\BaseNTok}[1]{\textcolor[rgb]{0.68,0.00,0.00}{#1}}
\newcommand{\BuiltInTok}[1]{\textcolor[rgb]{0.00,0.23,0.31}{#1}}
\newcommand{\CharTok}[1]{\textcolor[rgb]{0.13,0.47,0.30}{#1}}
\newcommand{\CommentTok}[1]{\textcolor[rgb]{0.37,0.37,0.37}{#1}}
\newcommand{\CommentVarTok}[1]{\textcolor[rgb]{0.37,0.37,0.37}{\textit{#1}}}
\newcommand{\ConstantTok}[1]{\textcolor[rgb]{0.56,0.35,0.01}{#1}}
\newcommand{\ControlFlowTok}[1]{\textcolor[rgb]{0.00,0.23,0.31}{\textbf{#1}}}
\newcommand{\DataTypeTok}[1]{\textcolor[rgb]{0.68,0.00,0.00}{#1}}
\newcommand{\DecValTok}[1]{\textcolor[rgb]{0.68,0.00,0.00}{#1}}
\newcommand{\DocumentationTok}[1]{\textcolor[rgb]{0.37,0.37,0.37}{\textit{#1}}}
\newcommand{\ErrorTok}[1]{\textcolor[rgb]{0.68,0.00,0.00}{#1}}
\newcommand{\ExtensionTok}[1]{\textcolor[rgb]{0.00,0.23,0.31}{#1}}
\newcommand{\FloatTok}[1]{\textcolor[rgb]{0.68,0.00,0.00}{#1}}
\newcommand{\FunctionTok}[1]{\textcolor[rgb]{0.28,0.35,0.67}{#1}}
\newcommand{\ImportTok}[1]{\textcolor[rgb]{0.00,0.46,0.62}{#1}}
\newcommand{\InformationTok}[1]{\textcolor[rgb]{0.37,0.37,0.37}{#1}}
\newcommand{\KeywordTok}[1]{\textcolor[rgb]{0.00,0.23,0.31}{\textbf{#1}}}
\newcommand{\NormalTok}[1]{\textcolor[rgb]{0.00,0.23,0.31}{#1}}
\newcommand{\OperatorTok}[1]{\textcolor[rgb]{0.37,0.37,0.37}{#1}}
\newcommand{\OtherTok}[1]{\textcolor[rgb]{0.00,0.23,0.31}{#1}}
\newcommand{\PreprocessorTok}[1]{\textcolor[rgb]{0.68,0.00,0.00}{#1}}
\newcommand{\RegionMarkerTok}[1]{\textcolor[rgb]{0.00,0.23,0.31}{#1}}
\newcommand{\SpecialCharTok}[1]{\textcolor[rgb]{0.37,0.37,0.37}{#1}}
\newcommand{\SpecialStringTok}[1]{\textcolor[rgb]{0.13,0.47,0.30}{#1}}
\newcommand{\StringTok}[1]{\textcolor[rgb]{0.13,0.47,0.30}{#1}}
\newcommand{\VariableTok}[1]{\textcolor[rgb]{0.07,0.07,0.07}{#1}}
\newcommand{\VerbatimStringTok}[1]{\textcolor[rgb]{0.13,0.47,0.30}{#1}}
\newcommand{\WarningTok}[1]{\textcolor[rgb]{0.37,0.37,0.37}{\textit{#1}}}

\usepackage{longtable,booktabs,array}
\newcounter{none} % for unnumbered tables
\usepackage{calc} % for calculating minipage widths
% Correct order of tables after \paragraph or \subparagraph
\usepackage{etoolbox}
\makeatletter
\patchcmd\longtable{\par}{\if@noskipsec\mbox{}\fi\par}{}{}
\makeatother
% Allow footnotes in longtable head/foot
\IfFileExists{footnotehyper.sty}{\usepackage{footnotehyper}}{\usepackage{footnote}}
\makesavenoteenv{longtable}
\usepackage{graphicx}
\makeatletter
\newsavebox\pandoc@box
\newcommand*\pandocbounded[1]{% scales image to fit in text height/width
  \sbox\pandoc@box{#1}%
  \Gscale@div\@tempa{\textheight}{\dimexpr\ht\pandoc@box+\dp\pandoc@box\relax}%
  \Gscale@div\@tempb{\linewidth}{\wd\pandoc@box}%
  \ifdim\@tempb\p@<\@tempa\p@\let\@tempa\@tempb\fi% select the smaller of both
  \ifdim\@tempa\p@<\p@\scalebox{\@tempa}{\usebox\pandoc@box}%
  \else\usebox{\pandoc@box}%
  \fi%
}
% Set default figure placement to htbp
\def\fps@figure{htbp}
\makeatother


% definitions for citeproc citations
\NewDocumentCommand\citeproctext{}{}
\NewDocumentCommand\citeproc{mm}{%
  \begingroup\def\citeproctext{#2}\cite{#1}\endgroup}
\makeatletter
 % allow citations to break across lines
 \let\@cite@ofmt\@firstofone
 % avoid brackets around text for \cite:
 \def\@biblabel#1{}
 \def\@cite#1#2{{#1\if@tempswa , #2\fi}}
\makeatother
\newlength{\cslhangindent}
\setlength{\cslhangindent}{1.5em}
\newlength{\csllabelwidth}
\setlength{\csllabelwidth}{3em}
\newenvironment{CSLReferences}[2] % #1 hanging-indent, #2 entry-spacing
 {\begin{list}{}{%
  \setlength{\itemindent}{0pt}
  \setlength{\leftmargin}{0pt}
  \setlength{\parsep}{0pt}
  % turn on hanging indent if param 1 is 1
  \ifodd #1
   \setlength{\leftmargin}{\cslhangindent}
   \setlength{\itemindent}{-1\cslhangindent}
  \fi
  % set entry spacing
  \setlength{\itemsep}{#2\baselineskip}}}
 {\end{list}}
\usepackage{calc}
\newcommand{\CSLBlock}[1]{\hfill\break\parbox[t]{\linewidth}{\strut\ignorespaces#1\strut}}
\newcommand{\CSLLeftMargin}[1]{\parbox[t]{\csllabelwidth}{\strut#1\strut}}
\newcommand{\CSLRightInline}[1]{\parbox[t]{\linewidth - \csllabelwidth}{\strut#1\strut}}
\newcommand{\CSLIndent}[1]{\hspace{\cslhangindent}#1}



\setlength{\emergencystretch}{3em} % prevent overfull lines

\providecommand{\tightlist}{%
  \setlength{\itemsep}{0pt}\setlength{\parskip}{0pt}}



 


% ---------------------------------------------------------------
% preamble.tex — ตั้งค่าภาษาไทยสำหรับ XeLaTeX
% ---------------------------------------------------------------

% 1) ตัดคำไทยด้วยพจนานุกรม ICU ภายใน XeTeX
%    (ภาษาไทยไม่มีช่องว่างระหว่างคำ ถ้าไม่ตั้งค่านี้บรรทัดจะไม่ตัดหรือล้นขอบ)
\XeTeXlinebreaklocale "th"
\XeTeXlinebreakskip = 0pt plus 1pt

% 2) ป้ายต่าง ๆ ของเล่มเป็นภาษาไทย
\AtBeginDocument{%
  \renewcommand{\contentsname}{สารบัญ}%
  \renewcommand{\listfigurename}{สารบัญรูป}%
  \renewcommand{\listtablename}{สารบัญตาราง}%
  \renewcommand{\bibname}{เอกสารอ้างอิง}%
  \renewcommand{\appendixname}{ภาคผนวก}%
  \renewcommand{\partname}{ภาค}%
  \renewcommand{\chaptername}{บทที่}%
}

% 3) หัวบทแบบ "บทที่ 1" ของ KOMA-Script
\renewcommand*{\chapterformat}{\chaptername~\thechapter\autodot\enskip}
\renewcommand*{\autodot}{}

% 4) ตัวหนาในหัวข้อ: ไม่ใช้ sans เพื่อให้ฟอนต์ไทยกลมกลืนทั้งเล่ม
\setkomafont{disposition}{\normalfont\bfseries}

% 5) ย่อหน้าบรรทัดแรก 1.5em และเว้นระยะเล็กน้อยระหว่างย่อหน้า
%    ปรับตัวเลขตามรสนิยมได้เลย
\setlength{\parindent}{1.5em}
\setlength{\parskip}{0.3em plus 0.1em}
\makeatletter
\@ifpackageloaded{tcolorbox}{}{\usepackage[skins,breakable]{tcolorbox}}
\@ifpackageloaded{fontawesome5}{}{\usepackage{fontawesome5}}
\definecolor{quarto-callout-color}{HTML}{909090}
\definecolor{quarto-callout-note-color}{HTML}{0758E5}
\definecolor{quarto-callout-important-color}{HTML}{CC1914}
\definecolor{quarto-callout-warning-color}{HTML}{EB9113}
\definecolor{quarto-callout-tip-color}{HTML}{00A047}
\definecolor{quarto-callout-caution-color}{HTML}{FC5300}
\definecolor{quarto-callout-color-frame}{HTML}{acacac}
\definecolor{quarto-callout-note-color-frame}{HTML}{4582ec}
\definecolor{quarto-callout-important-color-frame}{HTML}{d9534f}
\definecolor{quarto-callout-warning-color-frame}{HTML}{f0ad4e}
\definecolor{quarto-callout-tip-color-frame}{HTML}{02b875}
\definecolor{quarto-callout-caution-color-frame}{HTML}{fd7e14}
\makeatother
\makeatletter
\@ifpackageloaded{bookmark}{}{\usepackage{bookmark}}
\makeatother
\makeatletter
\@ifpackageloaded{caption}{}{\usepackage{caption}}
\AtBeginDocument{%
\ifdefined\contentsname
  \renewcommand*\contentsname{Table of contents}
\else
  \newcommand\contentsname{Table of contents}
\fi
\ifdefined\listfigurename
  \renewcommand*\listfigurename{List of Figures}
\else
  \newcommand\listfigurename{List of Figures}
\fi
\ifdefined\listtablename
  \renewcommand*\listtablename{List of Tables}
\else
  \newcommand\listtablename{List of Tables}
\fi
\ifdefined\figurename
  \renewcommand*\figurename{รูปที่}
\else
  \newcommand\figurename{รูปที่}
\fi
\ifdefined\tablename
  \renewcommand*\tablename{ตารางที่}
\else
  \newcommand\tablename{ตารางที่}
\fi
}
\@ifpackageloaded{float}{}{\usepackage{float}}
\floatstyle{ruled}
\@ifundefined{c@chapter}{\newfloat{codelisting}{h}{lop}}{\newfloat{codelisting}{h}{lop}[chapter]}
\floatname{codelisting}{Listing}
\newcommand*\listoflistings{\listof{codelisting}{List of Listings}}
\makeatother
\makeatletter
\makeatother
\makeatletter
\@ifpackageloaded{caption}{}{\usepackage{caption}}
\@ifpackageloaded{subcaption}{}{\usepackage{subcaption}}
\makeatother
\makeatletter
\@ifpackageloaded{sidenotes}{}{\usepackage{sidenotes}}
\@ifpackageloaded{marginnote}{}{\usepackage{marginnote}}
\makeatother
\usepackage{bookmark}
\IfFileExists{xurl.sty}{\usepackage{xurl}}{} % add URL line breaks if available
\urlstyle{same}
% fallback for those not using the hyperref driver hyperxmp:
\makeatletter
\@ifundefined{xmpquote}{\newcommand{\xmpquote}[1]{#1}}{}
\makeatother
\hypersetup{
  pdftitle={ระบาดวิทยาโรคข้าว},
  pdfauthor={กันต์ธณวิชญ์ ใจสงฆ์},
  colorlinks=true,
  linkcolor={Maroon},
  filecolor={Maroon},
  citecolor={Blue},
  urlcolor={Blue},
  pdfcreator={LaTeX via pandoc}}


\title{ระบาดวิทยาโรคข้าว}
\usepackage{etoolbox}
\makeatletter
\providecommand{\subtitle}[1]{% add subtitle to \maketitle
  \apptocmd{\@title}{\par {\large #1 \par}}{}{}
}
\makeatother
\subtitle{หลักการ การวิเคราะห์ และการประยุกต์ใช้}
\author{กันต์ธณวิชญ์ ใจสงฆ์}
\date{6 November 2026}
\begin{document}
\frontmatter
\maketitle

\renewcommand*\contentsname{Table of contents}
{
\setcounter{tocdepth}{1}
\tableofcontents
}

\setstretch{1.45}
\mainmatter
\bookmarksetup{startatroot}

\chapter*{คำนำ}\label{uxe04uxe33uxe19uxe33}
\addcontentsline{toc}{chapter}{คำนำ}

\markboth{คำนำ}{คำนำ}

ตำราเล่มนี้เขียนขึ้นสำหรับนักวิจัยและนักศึกษาที่ต้องการใช้หลักระบาดวิทยาพืช (plant disease
epidemiology) ในงานโรคข้าว เนื้อหาเริ่มจากแนวคิดเรื่องการเพิ่มขึ้นของโรคตามเวลา
ไปจนถึงการวิเคราะห์ข้อมูลจริงด้วยภาษา R

ทุกบทเขียนในรูปแบบ Quarto ตัวอย่างโค้ดจึงรันได้จริง
และผลลัพธ์ที่เห็นในเล่มมาจากโค้ดชุดเดียวกับที่ผู้อ่านคัดลอกไปใช้ได้

\bookmarksetup{startatroot}

\chapter{การเพิ่มขึ้นของโรคตามเวลา}\label{sec-disease-progress}

การเพิ่มขึ้นของโรคในแปลงข้าวไม่ได้เกิดแบบเส้นตรง ช่วงแรกโรคมักเพิ่มช้า
เพราะเนื้อเยื่อที่เป็นโรคยังมีน้อย จากนั้นจึงเร่งตัวขึ้นอย่างรวดเร็ว
และชะลอลงอีกครั้งเมื่อเนื้อเยื่อที่ยังไม่เป็นโรคเหลือน้อยลง รูปแบบนี้คือเส้นโค้งการเพิ่มขึ้นของโรค
(disease progress curve) ซึ่ง Van der Plank
(1963)\marginpar{\begin{footnotesize}
\begin{CSLReferences}{2}{0}
\bibitem[\citeproctext]{ref-vanderplank1963}
Van der Plank, J. E. 1963. \emph{Plant Diseases: Epidemics and Control}.
Academic Press.
\end{CSLReferences}
\vspace{2mm}\par\end{footnotesize}}
ใช้เป็นแกนหลักของระบาดวิทยาเชิงปริมาณ\sidenote{\footnotesize เชิงอรรถแบบนี้จะไปแสดงที่ขอบขวาของหน้าโดยอัตโนมัติ
  ทั้งในฉบับเว็บและฉบับ PDF}

\marginnote{\begin{footnotesize}

\textbf{ศัพท์ในบทนี้}\\
ความรุนแรงของโรค (disease severity)\\
อัตราการเพิ่มของโรค (apparent infection rate)\\
AUDPC

\end{footnotesize}}

\section{แบบจำลองพื้นฐาน}\label{sec-models}

ให้ \(y\) แทนสัดส่วนความรุนแรงของโรค (\(0 \le y \le 1\)) และ \(r\)
แทนอัตราการเพิ่มขึ้นของโรค (apparent infection rate) แบบจำลองลอจิสติก (logistic
model) เขียนในรูปสมการเชิงอนุพันธ์ได้ว่า

\begin{equation}\protect\phantomsection\label{eq-logistic}{
\frac{dy}{dt} = r\,y\,(1 - y)
}\end{equation}

เมื่ออินทิเกรตด้วยเงื่อนไขเริ่มต้น \(y(0) = y_0\) จะได้

\begin{equation}\protect\phantomsection\label{eq-logistic-sol}{
y(t) = \frac{1}{1 + \exp\!\left\{-\left[\ln\!\left(\dfrac{y_0}{1-y_0}\right) + r\,t\right]\right\}}
}\end{equation}

ส่วนแบบจำลองกอมเพิร์ตซ์ (Gompertz model) และแบบจำลองโมโนโมเลกุลาร์
(monomolecular model) มีรูปแบบดังนี้

\begin{equation}\protect\phantomsection\label{eq-gomp-mono}{
\frac{dy}{dt} = r\,y\,(-\ln y)
\qquad\text{และ}\qquad
\frac{dy}{dt} = r\,(1 - y)
}\end{equation}

Madden et al.
(2007)\marginpar{\begin{footnotesize}
\begin{CSLReferences}{2}{0}
\bibitem[\citeproctext]{ref-madden2007}
Madden, L. V., G. Hughes, and F. van den Bosch. 2007. \emph{The Study of
Plant Disease Epidemics}. APS Press.
\end{CSLReferences}
\vspace{2mm}\par\end{footnotesize}}
อธิบายการเลือกแบบจำลองและการประมาณค่าพารามิเตอร์ไว้อย่างละเอียด

\begin{tcolorbox}[enhanced jigsaw, arc=.35mm, bottomrule=.15mm, bottomtitle=1mm, breakable, colback=white, colbacktitle=quarto-callout-note-color!10!white, colframe=quarto-callout-note-color-frame, coltitle=black, left=2mm, leftrule=.75mm, opacityback=0, opacitybacktitle=0.6, rightrule=.15mm, title=\textcolor{quarto-callout-note-color}{\faInfo}\hspace{0.5em}{ข้อสังเกต}, titlerule=0mm, toprule=.15mm, toptitle=1mm]

แบบจำลองใน สมการ~\ref{eq-logistic} ให้อัตราการเพิ่มสูงสุดที่ \(y = 0.5\)
ขณะที่แบบกอมเพิร์ตซ์ให้จุดเปลี่ยนเว้า (inflection point) ที่
\(y = e^{-1} \approx 0.37\)

\end{tcolorbox}

\section{จำลองเส้นโค้งด้วย
R}\label{uxe08uxe33uxe25uxe2duxe07uxe40uxe2auxe19uxe42uxe04uxe07uxe14uxe27uxe22-r}

โค้ดต่อไปนี้วาดเส้นโค้งลอจิสติกจาก สมการ~\ref{eq-logistic-sol} ที่อัตรา \(r\)
ต่างกันสามระดับ ผลแสดงใน รูปที่~\ref{fig-logistic}

\begin{Shaded}
\begin{Highlighting}[]
\FunctionTok{library}\NormalTok{(tidyverse)}

\NormalTok{logistic }\OtherTok{\textless{}{-}} \ControlFlowTok{function}\NormalTok{(t, y0, r) }\DecValTok{1} \SpecialCharTok{/}\NormalTok{ (}\DecValTok{1} \SpecialCharTok{+} \FunctionTok{exp}\NormalTok{(}\SpecialCharTok{{-}}\NormalTok{(}\FunctionTok{qlogis}\NormalTok{(y0) }\SpecialCharTok{+}\NormalTok{ r }\SpecialCharTok{*}\NormalTok{ t)))}

\FunctionTok{expand\_grid}\NormalTok{(}\AttributeTok{t =} \FunctionTok{seq}\NormalTok{(}\DecValTok{0}\NormalTok{, }\DecValTok{60}\NormalTok{, }\AttributeTok{by =} \FloatTok{0.5}\NormalTok{), }\AttributeTok{r =} \FunctionTok{c}\NormalTok{(}\FloatTok{0.10}\NormalTok{, }\FloatTok{0.15}\NormalTok{, }\FloatTok{0.25}\NormalTok{)) }\SpecialCharTok{|\textgreater{}}
  \FunctionTok{mutate}\NormalTok{(}\AttributeTok{y =} \FunctionTok{logistic}\NormalTok{(t, }\AttributeTok{y0 =} \FloatTok{0.01}\NormalTok{, }\AttributeTok{r =}\NormalTok{ r),}
         \AttributeTok{r =} \FunctionTok{factor}\NormalTok{(r)) }\SpecialCharTok{|\textgreater{}}
  \FunctionTok{ggplot}\NormalTok{(}\FunctionTok{aes}\NormalTok{(t, y, }\AttributeTok{colour =}\NormalTok{ r)) }\SpecialCharTok{+}
  \FunctionTok{geom\_line}\NormalTok{(}\AttributeTok{linewidth =} \FloatTok{0.9}\NormalTok{) }\SpecialCharTok{+}
  \FunctionTok{labs}\NormalTok{(}\AttributeTok{x =} \StringTok{"วันหลังเริ่มพบโรค"}\NormalTok{, }\AttributeTok{y =} \StringTok{"ความรุนแรงของโรค (สัดส่วน)"}\NormalTok{,}
       \AttributeTok{colour =} \FunctionTok{expression}\NormalTok{(}\FunctionTok{italic}\NormalTok{(r))) }\SpecialCharTok{+}
  \FunctionTok{theme\_minimal}\NormalTok{(}\AttributeTok{base\_family =} \StringTok{"Sarabun"}\NormalTok{)}
\end{Highlighting}
\end{Shaded}

{
\makeatletter
\def\LT@makecaption#1#2#3{%
  \noalign{\smash{\hbox{\kern\textwidth\rlap{\kern\marginparsep
  \parbox[t]{\marginparwidth}{%
    \footnotesize{%
      \vspace{(1.1\baselineskip)}
    #1{#2: }\ignorespaces #3}}}}}}%
    }
\makeatother

\begin{figure}[H]

\sidecaption{\label{fig-logistic}เส้นโค้งลอจิสติกที่ \(y_0 = 0.01\)
และอัตราการเพิ่มของโรคสามระดับ}

\centering{

\pandocbounded{\includegraphics[keepaspectratio]{01-disease-progress_files/figure-pdf/fig-logistic-1.png}}

}

\end{figure}%

}

กราฟเล็กสำหรับเปรียบเทียบอาจวางไว้ที่ขอบขวาทั้งรูป เช่น
รูปที่~\ref{fig-margin-gompertz}

\begin{Shaded}
\begin{Highlighting}[]
\FunctionTok{tibble}\NormalTok{(}\AttributeTok{t =} \FunctionTok{seq}\NormalTok{(}\DecValTok{0}\NormalTok{, }\DecValTok{60}\NormalTok{, }\AttributeTok{by =} \FloatTok{0.5}\NormalTok{)) }\SpecialCharTok{|\textgreater{}}
  \FunctionTok{mutate}\NormalTok{(}\AttributeTok{y =} \FunctionTok{exp}\NormalTok{(}\FunctionTok{log}\NormalTok{(}\FloatTok{0.01}\NormalTok{) }\SpecialCharTok{*} \FunctionTok{exp}\NormalTok{(}\SpecialCharTok{{-}}\FloatTok{0.15} \SpecialCharTok{*}\NormalTok{ t))) }\SpecialCharTok{|\textgreater{}}
  \FunctionTok{ggplot}\NormalTok{(}\FunctionTok{aes}\NormalTok{(t, y)) }\SpecialCharTok{+}
  \FunctionTok{geom\_line}\NormalTok{(}\AttributeTok{linewidth =} \FloatTok{0.8}\NormalTok{) }\SpecialCharTok{+}
  \FunctionTok{labs}\NormalTok{(}\AttributeTok{x =} \StringTok{"วัน"}\NormalTok{, }\AttributeTok{y =} \StringTok{"y"}\NormalTok{) }\SpecialCharTok{+}
  \FunctionTok{theme\_minimal}\NormalTok{(}\AttributeTok{base\_family =} \StringTok{"Sarabun"}\NormalTok{, }\AttributeTok{base\_size =} \DecValTok{9}\NormalTok{)}
\end{Highlighting}
\end{Shaded}

\begin{marginfigure}

\sidecaption{\label{fig-margin-gompertz}เส้นโค้งกอมเพิร์ตซ์ที่ \(r = 0.15\)}

\centering{

\pandocbounded{\includegraphics[keepaspectratio]{01-disease-progress_files/figure-pdf/fig-margin-gompertz-1.png}}

}

\end{marginfigure}%

\section{พื้นที่ใต้เส้นโค้ง}\label{sec-audpc}

พื้นที่ใต้เส้นโค้งการเพิ่มขึ้นของโรค (area under the disease progress curve; AUDPC)
สรุปความรุนแรงตลอดฤดูเป็นตัวเลขเดียว โดยคำนวณด้วยวิธีสี่เหลี่ยมคางหมู (Shaner and
Finney
1977)\marginpar{\begin{footnotesize}
\begin{CSLReferences}{2}{0}
\bibitem[\citeproctext]{ref-shaner1977}
Shaner, G., and R. E. Finney. 1977. {``The Effect of Nitrogen
Fertilization on the Expression of Slow-Mildewing Resistance in {Knox}
Wheat.''} \emph{Phytopathology} 67: 1051--56.
\end{CSLReferences}
\vspace{2mm}\par\end{footnotesize}}

\begin{equation}\protect\phantomsection\label{eq-audpc}{
\text{AUDPC} = \sum_{i=1}^{n-1} \frac{y_i + y_{i+1}}{2}\,\left(t_{i+1} - t_i\right)
}\end{equation}

ตารางที่~\ref{tbl-audpc} เป็นข้อมูลสมมติเพื่อสาธิตการคำนวณเท่านั้น
ไม่ใช่ข้อมูลจากแปลงทดลองจริง

\begin{Shaded}
\begin{Highlighting}[]
\NormalTok{dat }\OtherTok{\textless{}{-}} \FunctionTok{tibble}\NormalTok{(}
  \AttributeTok{day   =} \FunctionTok{c}\NormalTok{(}\DecValTok{0}\NormalTok{, }\DecValTok{7}\NormalTok{, }\DecValTok{14}\NormalTok{, }\DecValTok{21}\NormalTok{, }\DecValTok{28}\NormalTok{),}
  \AttributeTok{var\_a =} \FunctionTok{c}\NormalTok{(}\DecValTok{1}\NormalTok{, }\DecValTok{4}\NormalTok{, }\DecValTok{12}\NormalTok{, }\DecValTok{30}\NormalTok{, }\DecValTok{45}\NormalTok{),}
  \AttributeTok{var\_b =} \FunctionTok{c}\NormalTok{(}\DecValTok{1}\NormalTok{, }\DecValTok{2}\NormalTok{, }\DecValTok{5}\NormalTok{, }\DecValTok{9}\NormalTok{, }\DecValTok{14}\NormalTok{)}
\NormalTok{)}

\NormalTok{audpc }\OtherTok{\textless{}{-}} \ControlFlowTok{function}\NormalTok{(t, y) }\FunctionTok{sum}\NormalTok{((}\FunctionTok{head}\NormalTok{(y, }\SpecialCharTok{{-}}\DecValTok{1}\NormalTok{) }\SpecialCharTok{+} \FunctionTok{tail}\NormalTok{(y, }\SpecialCharTok{{-}}\DecValTok{1}\NormalTok{)) }\SpecialCharTok{/} \DecValTok{2} \SpecialCharTok{*} \FunctionTok{diff}\NormalTok{(t))}

\NormalTok{dat }\SpecialCharTok{|\textgreater{}}
  \FunctionTok{add\_row}\NormalTok{(}\AttributeTok{day   =} \ConstantTok{NA}\NormalTok{,}
          \AttributeTok{var\_a =} \FunctionTok{audpc}\NormalTok{(dat}\SpecialCharTok{$}\NormalTok{day, dat}\SpecialCharTok{$}\NormalTok{var\_a),}
          \AttributeTok{var\_b =} \FunctionTok{audpc}\NormalTok{(dat}\SpecialCharTok{$}\NormalTok{day, dat}\SpecialCharTok{$}\NormalTok{var\_b)) }\SpecialCharTok{|\textgreater{}}
  \FunctionTok{mutate}\NormalTok{(}\AttributeTok{day =} \FunctionTok{if\_else}\NormalTok{(}\FunctionTok{is.na}\NormalTok{(day), }\StringTok{"AUDPC"}\NormalTok{, }\FunctionTok{as.character}\NormalTok{(day))) }\SpecialCharTok{|\textgreater{}}
\NormalTok{  knitr}\SpecialCharTok{::}\FunctionTok{kable}\NormalTok{(}\AttributeTok{align =} \StringTok{"lrr"}\NormalTok{, }\AttributeTok{col.names =} \FunctionTok{c}\NormalTok{(}\StringTok{"วัน"}\NormalTok{, }\StringTok{"พันธุ์ A"}\NormalTok{, }\StringTok{"พันธุ์ B"}\NormalTok{))}
\end{Highlighting}
\end{Shaded}

\begin{longtable}[]{@{}lrr@{}}

\caption{\label{tbl-audpc}ตัวอย่างสมมติ: ความรุนแรงของโรคไหม้ (\%) ในสองพันธุ์
และค่า AUDPC}

\tabularnewline

\toprule\noalign{}
วัน & พันธุ์ A & พันธุ์ B \\
\midrule\noalign{}
\endhead
\bottomrule\noalign{}
\endlastfoot
0 & 1 & 1.0 \\
7 & 4 & 2.0 \\
14 & 12 & 5.0 \\
21 & 30 & 9.0 \\
28 & 45 & 14.0 \\
AUDPC & 483 & 164.5 \\

\end{longtable}

จาก สมการ~\ref{eq-audpc} และ ตารางที่~\ref{tbl-audpc} จะเห็นว่าพันธุ์ B มี AUDPC
ต่ำกว่าพันธุ์ A อย่างชัดเจน ซึ่งเป็นลักษณะของความต้านทานแบบลดอัตรา (rate-reducing
resistance) ที่อธิบายไว้ใน หัวข้อ~\ref{sec-models}

\bookmarksetup{startatroot}

\chapter*{เอกสารอ้างอิง}\label{uxe40uxe2duxe01uxe2auxe32uxe23uxe2duxe32uxe07uxe2duxe07}
\addcontentsline{toc}{chapter}{เอกสารอ้างอิง}

\markboth{เอกสารอ้างอิง}{เอกสารอ้างอิง}



\backmatter


\end{document}
